\sectioncentered*{Реферат}

РАЗРАБОТКА СПЕЦИФИКАЦИИ НА ОСНОВЕ ПРАВИЛ И РЕАЛИЗАЦИЯ ПРОГРАММЫ ДЛЯ ДИАГНОСТИКИ ПРИЧИН НЕИСПРАВНОСТИ АВТОМОБИЛЯ: курсовой проект /С.В. Марковцев.. – Минск : БГУИР, 2025, – п.з. – 41 с., чертежей (плакатов) – 1 л. формата А1. 

Курсовой проект на тему "Разработка спецификации на основе правил и реализация программы для диагностики причин неисправности автомобиля" разработан с целью повышения степени удовлетворенности клиентов автосервиса за счет анализа накопленной информации о клиентском поведении, оптимизации работы механика.

Пояснительная записка к курсовому проекту состоит из введения, 3 разделов, включающих, анализ технического задания и создания спецификации к нему, проектирования системы, а именно создание различных диаграмм, реализацию системы, заключение, список использованных источников и приложение, содержащее листинг кода отдельных классов.

Для разработки проекта системы выбран объектно-ориентированный язык проектирования \textit{UML}. Проект представляет собой описание создания библиотечной системы. В нем присутствуют различные диаграммы, каждая из которых демонстрирует определенный аспект проекта.

В результате выполнения курсовой работы спроектирована модель системы, написан глоссарий, а также составлены чертежи. 

Результаты, полученные в ходе курсового проектирования, могут использоваться участниками различных проектов, где требуется работа с обширной клиентской базой. Разработка спецификации на основе правил и реализация программы для диагностики причин неисправности автомобиля является важным процессом, позволяющим автоматизировать и упростить процесс определения неисправностей. Разработанная спецификация и программная реализация могут значительно сократить время и усилия, затрачиваемые на диагностику автомобиля, и повысить эффективность устранения неисправностей. Эта модель имеет необходимые диаграммы для написания ядра системы, которая обрабатывает запросы и помогает автоматизировать деятельность работника автосервиса. Созданная модель является также хорошим примером создания \textit{UML} диаграмм и может служить примером написания моделей других проектов. 